% Generated draft: requires mathematical and source review.
\documentclass[11pt]{article}
\usepackage[margin=25mm]{geometry}
\usepackage{fontspec}
\setmainfont{Latin Modern Roman}
\usepackage{amsmath,amssymb,mathtools,unicode-math}
\setmathfont{Latin Modern Math}
\usepackage{enumitem,needspace,xurl,hyperref}
\setlength{\parindent}{0pt}
\setlength{\parskip}{5pt}
\setlength{\emergencystretch}{3em}
\newcommand{\N}{\mathbb{N}}
\newcommand{\Z}{\mathbb{Z}}
\newcommand{\Q}{\mathbb{Q}}
\newcommand{\R}{\mathbb{R}}
\newcommand{\C}{\mathbb{C}}
\newcommand{\F}{\mathbb{F}}
\newcommand{\A}{\mathbb{A}}
\newcommand{\PP}{\mathbb{P}}
\newcommand{\E}{\mathbb{E}}
\newcommand{\eps}{\varepsilon}
\newcommand{\dd}{\,\mathrm{d}}
\begin{document}
\section*{UnsolvedMath ID: 30002201}
\textbf{Source problem code:} catalogue entry
\subsection{Definitions and mathematical statement}
The supplied record concerns a wonderful compactification \(Y\) of a hyperplane-arrangement complement. It does not specify the ground field, the arrangement, the building set defining the wonderful compactification, or the meaning of the referenced tropical iterative procedure.

For clarity, one possible convention for an effective cone can be described conditionally. If \(Y\) is a projective variety over a field, let \(N^1(Y)_{\R}\) denote the real vector space of Cartier divisors modulo numerical equivalence, where numerical equivalence means equality of intersection degrees with every integral complete curve in \(Y\). Under the effective Cartier divisor convention, set
\[
\operatorname{Eff}(Y)=\left\{\sum_{i=1}^{r} a_i[D_i] : r\in\Z_{\geq 0},\ a_i\in\R_{\geq 0},\ D_i\text{ an effective Cartier divisor on }Y\right\}\subseteq N^1(Y)_{\R}.
\]
Its closure \(\overline{\operatorname{Eff}(Y)}\) is the corresponding pseudoeffective cone. These definitions explain a possible interpretation only: the source does not specify the divisor convention, ambient vector space, or whether closure is intended.

\textbf{Source question, with unresolved qualifications.} For a wonderful compactification \(Y\) of a hyperplane-arrangement complement, can the tropical iterative procedure determine its effective cone and provide a method for studying the Castravet--Tevelev conjecture referred to in the source as concerning \(M_{0,n}\)?

The source does not define the procedure, state its inputs or output, assert termination, or specify whether the question concerns every such \(Y\) or a restricted class. It also does not state the Castravet--Tevelev conjecture or specify the range of \(n\). The notation \(M_{0,n}\) ordinarily refers to the moduli space of smooth genus-zero curves with \(n\) ordered distinct marked points, whereas \(\overline{M}_{0,n}\) denotes its stable compactification; the intended space here is unresolved. Consequently, the supplied question is a research direction rather than a fully specified mathematical proposition. No replacement of \(M_{0,n}\) by its compactification, or full conjectural assertion, is inferred.
\subsection{Short English statement}
The question asks whether an unspecified tropical iterative procedure can compute effective cones of wonderful arrangement compactifications and help study a Castravet--Tevelev conjecture. The procedure, the class of compactifications, the effective-cone convention, and the intended moduli space and conjecture must be clarified before a precise theorem or algorithmic claim can be stated.
\subsection{Sources}
\begin{itemize}
\item Supplied local UnsolvedMath registry record: permanent ID \(30002201\), title \emph{Effective Cones of Wonderful Arrangement Compactifications}, source problem code \texttt{catalogue entry}, in the Oberwolfach Reports Open Problems collection. No individual proposer is supplied; the proposed year is recorded as 2012.
\item Supplied report citation: \emph{Mini-Workshop: Cohomology Rings and Fundamental Groups of Hyperplane Arrangements, Wonderful Compactifications, and Real Toric Varieties} (2013), Oberwolfach Reports, pp.~2939--2983, DOI \texttt{10.4171/owr/2012/49}. Supplied source URL and literature-source link, labeled ``original OWR report'': \url{https://doi.org/10.4171/owr/2012/49}. The background describes these reports as documenting mathematical workshops and their open-problem discussions.
\item Literature assessment dated \texttt{2026-08-21}: the supplied curation labels the problem open and says that the extracted record only suggests that an iterative procedure may help study the Castravet--Tevelev effective-cone conjecture; it neither states that conjecture in full nor yields a determinate resolution. This is supplied dated curation, not an independently verified literature assessment. The source expressly requests independent verification of time-sensitive status and cited claims.
\item The registry supplies statement status \texttt{corrected\_verified} and reports that the mathematical content was checked against the official Oberwolfach report, with notation, context, and question form normalized without changing the source-level claim. This verification assertion is supplied provenance and has not been independently checked for this draft.
\end{itemize}
\end{document}
